\chapter{Velké jazykové model}
\label{ch:llm}

Tato kapitola se věnuje problematikou velkých jazykových modelů (LLM) od jejich historických kořenů přes architektonické jádro -- Transformer -- až po způsob trénování, fine-tuningu a nasazení.
Cílem je poskytnout dostatečný teoretický základ pro pochopení designových rozhodnutí popsaných v následujících kapitolách.

Druhá část kapitoly zasazuje LLM do kontextu analýzy zdrojového kódu: co je k takové analýze potřeba, jakým způsobem ji model provádí a v čem se od tradičních statických nástrojů liší.
Kapitola uzavírá přehledem hlavních rizik a omezení, která je nutné zohlednit při nasazení v reálném školním prostředí.

% Tato kapitola se věnuje velkým jazykovým modelům (LLM) z pohledu jejich principu fungování a následně je zasazuje do kontextu analýzy zdrojového kódu.
% V závěru kapitoly jsou LLM porovnány s tradičními statickými analyzátory a jsou rozebrána hlavní rizika a omezení, která je nutné zohlednit při jejich nasazení do školního systému.

\input{chapters/llm/3-1-llm-principes}
\section{Analýza zdrojového kódu pomocí LLM}
\label{sec:llm-analyze}

\TODO{Stále nevím, zda celou tuto sekci nebudu psát v implementaci. Proto zatím nepíšu.}

\endinput


\input{chapters/llm/3-3-llm-agains-static}
\input{chapters/llm/3-4-llm-limitations}

\endinput